\chapter{Security Analysis}

\section{Security of AES-CTR-HMAC with RISC Zero Proofs}

\subsection{Scope}

This chapter gives a game-based proof sketch, not a full formal reduction. The
goal is to explain why the public transcript produced by the implemented
pipeline does not reveal which of two equal-length plaintexts was encrypted,
assuming the usual cryptographic and zkVM assumptions hold.

The public challenge transcript is written as
\((\mathrm{aad}, n, c, t, \pi)\), where \(\mathrm{aad}\) is associated data or
context, \(n\) is the AES-CTR IV/nonce, \(c\) is ciphertext, \(t\) is the
truncated HMAC tag, and \(\pi\) is the RISC Zero receipt. If a future guest also
journals application output \(y_{pub}\), confidentiality requires that
\(y_{pub}\) be independent of the challenge bit except for intentional public
metadata such as message length. A plaintext-dependent journal value can reveal
information regardless of the encryption layer.

The symbol \(Y\) below denotes the internal AES-CTR keystream: the output of
AES-128 evaluated on the nonce-derived counter blocks and truncated to the
message length. It is not a public journal value.

The implemented MAC input is also explicit. It is not a bare concatenation of
fields. Let \(D\) be the domain separator
\texttt{aes-ctr+hmac-sha256-192:v1}, and let \(|x|_{64}\) denote the 64-bit
big-endian encoding of a byte-string length. The framed byte string is:

\[
\mathsf{Frame}(\mathrm{aad}, n, c) =
D \parallel |\mathrm{aad}|_{64} \parallel \mathrm{aad} \parallel
|n|_{64} \parallel n \parallel |c|_{64} \parallel c.
\]

This matches the shared \texttt{compute\_auth\_tag\_192(...)} helper used by the
host and guest code.

This chapter states the confidentiality, integrity, and execution-binding
assumptions used by the report, but it does not prove full chosen-ciphertext
security, sender identity, replay protection, or side-channel resistance.

\subsection{Assumptions}

The proof sketch relies on the following assumptions:

\begin{itemize}
\item \textbf{Nonce discipline:} AES-CTR IV/nonce values are unique under each
encryption key. Equivalently, the same AES counter block is never encrypted twice
under the same key, including across multi-block messages and counter increments.
If two messages reuse the same counter stream, XORing their ciphertexts cancels
the keystream and leaks the XOR of the plaintexts.
\item \textbf{AES-128-CTR confidentiality:} AES-128 is treated as a secure block
cipher/PRF, so CTR mode is IND-CPA secure when nonces are never reused
\cite{nistfips197,nistsp80038a,katz2020introduction}.
\item \textbf{HMAC security:} HMAC-SHA256, truncated to 192 bits in this
implementation, is treated as a secure PRF for tag indistinguishability and as a
secure MAC for tamper-detection reasoning
\cite{rfc2104,nistfips1981,nistfips1804}.
\item \textbf{Key separation:} encryption and MAC keys are distinct.
\item \textbf{zkVM proof properties:} receipt verification is sound for the
expected guest image ID, and the proof does not leak non-journaled witness data
beyond the public journal and verification result.
\end{itemize}

The 192-bit tag truncation gives a generic blind-forgery term of approximately
\(q_v / 2^{192}\) for \(q_v\) verification attempts, before accounting for any
advantage against HMAC itself. This is negligible for the scale considered in
this dissertation, but it is still a stated cryptographic assumption rather than
an empirical benchmark result.

\subsection{IND-CPA Security Game with Encryption Oracle}

The proof-output confidentiality game is an IND-CPA-style game with an explicit
encryption oracle. It models the fact that an adversary may observe other valid
authenticated proof transcripts before and after the challenge transcript.

\begin{enumerate}
\item The challenger samples independent keys \(k_E\) and \(k_M\), initializes an
empty nonce set \(S\), and samples the challenge bit \(b \leftarrow \{0,1\}\).
\item The adversary \(\mathcal{A}\) receives access to an encryption oracle
\(\mathcal{O}_{enc}\). On input \((m, \mathrm{aad})\), the oracle samples a fresh
nonce \(n \notin S\), adds it to \(S\), computes the AES-CTR keystream
\(Y = \mathrm{AES\text{-}CTR\text{-}Stream}_{k_E}(n, |m|)\), sets
\(c = m \oplus Y\), and computes
\[
t = \mathrm{Trunc}_{192}(\mathrm{HMAC}_{k_M}(\mathsf{Frame}(\mathrm{aad}, n, c))).
\]
It then proves the guest relation for the public journal
\((\mathrm{aad}, n, c, t)\) and returns \((\mathrm{aad}, n, c, t, \pi)\).
\item At challenge time, \(\mathcal{A}\) outputs two equal-length messages
\(m_0, m_1\) and challenge associated data \(\mathrm{aad}^*\).
\item The challenger samples a fresh nonce \(n^* \notin S\), adds it to \(S\),
computes
\[
Y^* = \mathrm{AES\text{-}CTR\text{-}Stream}_{k_E}(n^*, |m_b|),
\quad c^* = m_b \oplus Y^*,
\]
and computes
\[
t^* = \mathrm{Trunc}_{192}(\mathrm{HMAC}_{k_M}(\mathsf{Frame}(\mathrm{aad}^*, n^*, c^*))).
\]
The challenger returns \((\mathrm{aad}^*, n^*, c^*, t^*, \pi^*)\), where
\(\pi^*\) is a receipt for the expected guest image and journal.
\item \(\mathcal{A}\) may continue to query \(\mathcal{O}_{enc}\), with the same
fresh-nonce rule, and finally outputs a guess \(b'\).
\end{enumerate}

The advantage is:

\[
\mathrm{Adv}_{\mathcal{A}} :=
\left| \Pr[b' = b] - \tfrac{1}{2} \right|.
\]

If the oracle is restricted to the single challenge query only, this game reduces
to the simpler one-shot setting. The oracle version is stated here because it is
the standard IND-CPA shape and makes nonce uniqueness explicit across multiple
encryptions.

\subsection{Game Sequence}

\paragraph{Game \(G_0\) (real execution).}
This is the game above, using AES-128-CTR, HMAC-SHA256-192 over
\(\mathsf{Frame}\), and real RISC Zero receipts.

\paragraph{Game \(G_1\) (replace AES-CTR keystreams).}
For every oracle response and for the challenge response, replace the internal
AES-CTR keystream \(Y\) with a uniformly random string \(R\) of the same length:

\[
c = m \oplus R, \quad R \leftarrow \{0,1\}^{|m|}.
\]

Here \(Y\) is the concatenation of AES-128 outputs on the nonce-derived counter
blocks \((n, n+1, \ldots)\), truncated to the plaintext length. Under the
IND-CPA security of nonce-respecting AES-CTR:

\[
|\Pr[G_0] - \Pr[G_1]| \leq
\mathrm{Adv}_{\mathrm{IND\text{-}CPA}}^{\mathrm{AES\text{-}CTR}}(\mathcal{A}).
\]

Because the challenge messages have equal length, \(c^* = m_b \oplus R^*\) is
distributed independently of \(b\) in \(G_1\).

\paragraph{Game \(G_2\) (replace HMAC tags).}
Replace each HMAC tag on \(\mathsf{Frame}(\mathrm{aad}, n, c)\) with a random
192-bit value. After \(G_1\), the frame contains only public values whose
challenge ciphertext component is independent of \(b\). This hop records the
assumption that truncated HMAC output does not introduce a distinguishable signal
about \(b\):

\[
|\Pr[G_1] - \Pr[G_2]| \leq
\mathrm{Adv}_{\mathrm{PRF}}^{\mathrm{HMAC}}(\mathcal{A}).
\]

For integrity rather than confidentiality, the corresponding forgery term is
bounded informally by the HMAC unforgeability advantage plus the truncation term
\(q_v / 2^{192}\).

\paragraph{Game \(G_3\) (simulate receipts).}
Let \(\mathrm{Sim}\) be the zero-knowledge simulator for the proof system. Given a
public journal \(J = (\mathrm{aad}, n, c, t)\) and the expected guest image, it
outputs a simulated receipt \(\tilde{\pi}\) that is computationally
indistinguishable from a real receipt for the same public journal. Replace each
real receipt \(\pi\) with \(\tilde{\pi} \leftarrow \mathrm{Sim}(J)\).

By the zero-knowledge property for non-journaled witness data:

\[
|\Pr[G_2] - \Pr[G_3]| \leq
\mathrm{Adv}_{\mathrm{ZK}}(\mathcal{A}).
\]

This step does not hide the journal. It only says the receipt should not reveal
extra information about secret inputs beyond the values intentionally committed
to the journal.

\paragraph{Game \(G_4\) (ideal).}
In the final game, the challenge tuple \((c^*, t^*, \pi^*)\) is independent of
\(b\). Therefore:

\[
\Pr[G_4 : b' = b] = \tfrac{1}{2}.
\]

\subsection{Final Bound}

By the triangle inequality, the confidentiality advantage is bounded by:

\begin{align*}
\mathrm{Adv}_{\mathcal{A}}
&\leq \mathrm{Adv}_{\mathrm{IND\text{-}CPA}}^{\mathrm{AES\text{-}CTR}}
+ \mathrm{Adv}_{\mathrm{PRF}}^{\mathrm{HMAC}}
+ \mathrm{Adv}_{\mathrm{ZK}}.
\end{align*}

The AES-CTR term covers the encryption-oracle queries under the nonce-uniqueness
assumption. The HMAC term covers tag indistinguishability for confidentiality;
integrity relies instead on MAC unforgeability plus the blind-forgery term
\(q_v / 2^{192}\). Execution binding relies separately on RISC Zero soundness and
verification against the expected image ID.

\subsection{Framing and Integrity Note}

Length-prefixing the authenticated frame is required for integrity. Without it,
the byte string \(\mathrm{aad} \parallel n \parallel c\) can be ambiguous when
variable-length AAD is allowed: an adversary could shift bytes between AAD and
ciphertext while preserving the same concatenated MAC input. The implemented
frame prevents this by authenticating a domain separator and every field length
before the corresponding bytes. This binds field boundaries, so a valid tag for
one tuple cannot be reinterpreted as a valid tag for a different tuple with a
different AAD/ciphertext split.

\subsection{Conclusion}

Under nonce-respecting AES-128-CTR, separated AES and HMAC keys, PRF/MAC security
of HMAC-SHA256-192, length-prefixed authenticated framing, and RISC Zero proof
soundness and non-journal leakage assumptions, the public transcript does not
reveal which equal-length challenge plaintext was encrypted. The result is
deliberately scoped: it supports the report's confidentiality and
transcript-binding claims for the implemented benchmark pipeline, but it is not a
full AEAD proof or a deployment-level network security proof.

\hfill \(\rule{0.8ex}{0.8ex}\)
